\hspace{3mm}

\centerline{\bf \Huge Домашнее задание.}
\centerline{\bf \Huge Принцип \href{https://vk.com/video-91031095_456245511}{Дирихле}}

\begin{flushright}
    \scriptsize{Насладись небольшим отклонением от пути.\\ По полной.\\ Там ты сможешь найти то, что важнее желаемого тобой. \\Джин Фрикс. Хантер}
\end{flushright}

\problem{1!} Верно ли, что среди любых 10 целых чисел можно выбрать хотя бы две пары чисел(одно число в разных парах может повторяться), разность которых кратна 8?

\problem{2!} В команде ЭлМШ 53 человека:
Шиншиллы - 22 человека, Лососи - 13, Киви - 14, АА - 1, РВ - 1, ЕМ - 1, ВГ(контент менеджер) - 1 человек. Какое наименьшее число людей нужно выбрать, чтобы среди них заведомо оказалось не менее 14 людей одного типа?

\problem{3}Кирилл провел на плоскости 36 произвольных прямых. Докажите, что найдутся две прямые угол между которыми не более чем 5 градусов

\problem{4}Очир Владимирович выписал на доску 100 различных целых чисел. Подошел Анджур Аркадьевич и выбрал из них несколько чисел так, что их сумма делится на 100. Очиру Владимировичу это не понравилось и он стер все числа с доски и написал новые. Докажите, что какие бы числа не выписал Очир Владимирович на доску, подоспевшему Расулу Владимировичу всё равно удастся найти несколько чисел таких, что их сумма делится на 100.